\chapter{Framen}
\label{ch:frameformat}

{\large\itshape Fält för fält, från SOF till EOF.\par}

\begin{chapterbox}{I det här kapitlet}
\begin{itemize}
  \item Dataframens sju fält, och vad vart och ett är till för.
  \item Identifieraren, som är två saker samtidigt: adress och prioritet.
  \item Varför CAN inte har någon destinationsadress.
  \item DLC, och varför den är fyra bitar när nyttolasten är högst åtta byte.
  \item Kvittensen som alla mottagare skriver i samma bit.
\end{itemize}
\end{chapterbox}

\section{Framens form}

En CAN-dataframe med 11-bitars identifierare ser ut så här:

\bookfigure{frame_layout}{Dataframens fält. Bredderna är i bitar, före bitstoppning.}{fig:frame}

Sju fält, och en gräns som är viktigare än de andra: allt fram till och med RTR-biten är
\textbf{arbitreringsfältet}, och det är där flera samtidigt sändande noder avgör vem som får
fortsätta (\kapref{ch:arbitrering}).

\begin{booktable}{@{}llL@{}}
\thead{Fält} & \thead{Bitar} & \thead{Innehåll} \\ \midrule
SOF & 1 & Start of frame. En dominant bit, som säger att en frame börjar. \\
Identifierare & 11 & Vad framen innehåller, och därmed dess prioritet. \\
RTR & 1 & Remote transmission request: dominant för en dataframe. \\
IDE & 1 & Identifier extension: dominant för 11-bitars identifierare. \\
r0 & 1 & Reserverad, dominant. \\
DLC & 4 & Data length code: antal databytes, 0--8. \\
Data & 0--64 & Nyttolasten, noll till åtta byte. \\
CRC & 15 & Checksumma över allt föregående. \\
CRC-avgränsare & 1 & Recessiv. \\
ACK-lucka & 1 & Sändaren släpper; en mottagare drar ned. \\
ACK-avgränsare & 1 & Recessiv. \\
EOF & 7 & End of frame, sju recessiva bitar. \\
\end{booktable}

Efter EOF följer dessutom \textbf{intermission}, tre recessiva bitar då ingen får börja sända.
Tillsammans med EOF ger det tio recessiva bitar i rad mellan två frames --- ett mönster som inte
kan uppstå någon annanstans, eftersom bitstoppningen (\kapref{ch:stuffing}) gör att fem lika
bitar i rad är det mesta som förekommer inuti en frame.

\begin{aside}
\note[Räkna] $1 + 12 + 6 + 64 + 16 + 2 + 7 = 108$ bitar för den längsta dataframen, före
stoppning. Bitstoppningen lägger till som mest ett tjugotal, så en maximal frame är omkring
130 bitar. Vid 1~Mbit/s tar den alltså ungefär 130~µs. Den siffran är värd att komma ihåg: den
är svaret på "hur länge kan en sändning hålla på?".
\end{aside}

\section{Identifieraren}

Identifieraren är CAN:s mest särpräglade designval, och den är två saker samtidigt.

\subsection{Den säger vad framen innehåller}

Den är \textbf{inte} en mottagaradress. Den beskriver \emph{innehållet}:

\begin{narrowtable}{@{}ll@{}}
\thead{Identifierare} & \thead{Betydelse} \\ \midrule
\code{0x301} & Motorvarvtal \\
\code{0x302} & Kylvätsketemperatur \\
\code{0x400} & Statusförfrågan \\
\end{narrowtable}

En nod som sänder \code{0x301} säger inte "till instrumentpanelen", den säger "varvtalet är
detta". Vem som lyssnar är inte sändarens problem, och det får konsekvenser som är lätta att
missa i början:

\begin{itemize}
  \item \textbf{Flera mottagare kan reagera på samma frame.} Både instrumentpanelen och
    växellådsstyrningen kan behöva varvtalet, och ingen av dem behöver veta om den andra.
  \item \textbf{En ny nod kan läggas till utan att någon gammal ändras.} Den börjar bara lyssna
    efter de identifierare den bryr sig om.
  \item \textbf{Filtreringen ligger hos mottagaren.} Varje nod tar emot varje frame och avgör
    själv om den angår den --- i mjukvara, eller i kontrollerns filterhårdvara.
\end{itemize}

\subsection{Den bestämmer prioritet}

Samtidigt är identifieraren framens prioritet, och där gäller att \textbf{lägre numeriskt värde
vinner}. Det följer direkt av \kapref{ch:buss}: identifieraren sänds MSB först, en nolla är
dominant, och den dominanta biten vinner.

Det betyder att adressrymden måste planeras. En identifierare är inte bara ett namn; den är en
plats i en prioritetsordning som gäller hela bussen. Bromsar får låga nummer, kupévärme höga.

\begin{aside}
\note[En vanlig fallgrop] Två noder får aldrig sända \emph{samma} identifierare. Två frames med
samma identifierare är bit för bit identiska genom hela arbitreringsfältet, så ingen av noderna
upptäcker den andra --- och så fort deras datafält skiljer sig åt uppstår ett bitfel mitt i
framen. Identifierare tilldelas per sändande nod, inte per meddelandetyp.
\end{aside}

\subsection{Utökade identifierare}

CAN 2.0B tillåter 29-bitars identifierare. Skillnaden syns i IDE-biten: recessiv i stället för
dominant, och därefter följer ytterligare 18 identifierarbitar. Den här boken håller sig till
11 bitar, men det är värt att veta varför valet syns i framen: en nod som bara förstår 11 bitar
måste kunna se, tidigt och entydigt, att den inte förstår den här framen.

Notera också prioritetskonsekvensen: en 11-bitars frame med samma inledande bitar vinner alltid
över en 29-bitars, eftersom dess IDE-bit är dominant.

\section{RTR och fjärrframes}

RTR-biten skiljer två sorters frames:

\begin{narrowtable}{@{}lll@{}}
\thead{RTR} & \thead{Sorts frame} & \thead{Betydelse} \\ \midrule
\dom{} (0) & Dataframe & "Här är datan." \\
\rec{} (1) & Fjärrframe & "Skicka mig datan." \\
\end{narrowtable}

En \textbf{fjärrframe} har samma identifierare som den dataframe den efterfrågar, en DLC men
\emph{inget} datafält. Tanken är att en nod ska kunna be om ett värde i stället för att vänta på
att det sänds.

I praktiken används fjärrframes sällan i moderna konstruktioner. De komplicerar mer än de
tillför: två noder kan svara på samma begäran, och en nod som inte förstår mekanismen ser en
frame med en DLC som inte stämmer med datafältets längd.

Att RTR ligger \emph{i} arbitreringsfältet är däremot inte godtyckligt. Det gör att en dataframe
alltid vinner över en fjärrframe med samma identifierare --- den som har datan får företräde över
den som ber om den, vilket är precis rätt prioritering.

\section{Kontrollfältet och DLC}

Tre bitar administration, sedan fyra bitar längd:

\begin{itemize}
  \item \textbf{IDE} --- dominant för 11 bitars identifierare.
  \item \textbf{r0} --- reserverad, sänds dominant. Den finns för framtida bruk, och den
    framtiden kom aldrig för klassisk CAN.
  \item \textbf{DLC} --- fyra bitar, värde 0--8.
\end{itemize}

Fyra bitar rymmer 0--15, men bara 0--8 är meningsfulla. Vad händer med 9--15? I standarden är de
reserverade och ska behandlas som 8. I praktiken är det en av de punkter där kontrollrar skiljer
sig åt --- och det är därför en drivrutin som \emph{sänder} aldrig ska producera en DLC över 8,
oavsett vad hårdvaran råkar göra med den.

\section{CRC-fältet}

Femton bitars checksumma, beräknad över allt från SOF till och med datafältet, följd av en
recessiv avgränsare. Hur den beräknas står i \kapref{ch:stuffing}.

Avgränsaren finns av en praktisk anledning: den ger mottagarna en känd, recessiv bit att ta
sats i innan ACK-luckan, som är den enda bit i hela framen där någon annan än sändaren driver
bussen.

\section{ACK-fältet}

Här händer något som inte har någon motsvarighet i ett hemsnickrat protokoll.

Sändaren sänder ACK-luckan \textbf{recessivt} --- alltså släpper den. Varje mottagare som har
tagit emot framen och fått CRC:n att stämma drar samtidigt ned den \textbf{dominant}. Sändaren
läser tillbaka bussen och ser:

\begin{narrowtable}{@{}ll@{}}
\thead{Sändaren läser} & \thead{Slutsats} \\ \midrule
\dom{} & Minst en nod tog emot framen felfritt \\
\rec{} & Ingen gjorde det \\
\end{narrowtable}

Tre saker är värda att stanna vid:

\textbf{Kvittensen kostar ingenting.} Den är en enda bit, inuti framen, och den kräver ingen
svarsframe. Jämför med ett ACK-meddelande som måste konstrueras, sändas, arbitreras om och
kvitteras i sin tur.

\textbf{Den säger inte \emph{vem}.} Alla mottagare drar ned samma bit samtidigt. Sändaren vet att
\emph{någon} hörde, aldrig hur många eller vilka. Behövs det måste applikationen svara med en
egen frame.

\textbf{Den säger inte att någon \emph{brydde sig}.} En nod som filtrerar bort framen som
ointressant kvitterar den ändå, för den tog emot den felfritt. Kvittensen är ett kvitto på
\emph{överföringen}, inte på \emph{behandlingen}.

\begin{simplification}
\note[Förenkling] Den förenklade kontrollern i \kapref{ch:kursen} \emph{släpper} ACK-luckan som
sändare och \emph{drar ned} den som mottagare, precis enligt standarden --- men den läser aldrig
tillbaka vad som hände där. En frame som ingen kvitterar fullbordas därför som om allt gått bra.
\end{simplification}

\section{EOF och intermission}

Sju recessiva bitar avslutar framen, och tre till av intermission håller bussen ledig innan
någon får börja om.

Sju recessiva i rad är fler än bitstoppningen någonsin tillåter inuti en frame, vilket gör EOF
till ett entydigt slutmärke: en mottagare som ser sju recessiva bitar vet att framen är slut,
oavsett var den tappade tråden.

Det är också här felsignalering sker: en nod som upptäcker ett fel avbryter genom att sända sex
dominanta bitar i rad --- ett mönster som är omöjligt i en korrekt frame, och som därför inte kan
förväxlas med data. \Kapref{ch:fel} går igenom det.

\section*{Övningar}

\exercise{Läs framen}{Avkodning}
\label{ex:2:1}
En frame har identifieraren \code{0x2A7}, DLC 5 och databytesen \code{01 02 03 04 05}.
\begin{enumerate}[a)]
  \item Hur många bitar är framen, före bitstoppning?
  \item Hur lång tid tar den vid 1~Mbit/s, före stoppning?
  \item Vilken identifierare skulle ha vunnit över den i en arbitrering: \code{0x2A6} eller
    \code{0x2A8}?
\end{enumerate}

\exercise{Ingen adress}{Resonemang}
\label{ex:2:2}
Ett system har en temperatursensor och två noder som båda behöver temperaturen.
\begin{enumerate}[a)]
  \item Hur många frames måste sensorn sända?
  \item Hur skulle samma sak sett ut i ett protokoll med destinationsadress?
  \item En tredje nod tillkommer som också vill ha temperaturen. Vad måste ändras i sensorn?
\end{enumerate}

\exercise{Kvittensen}{Resonemang}
\label{ex:2:3}
En nod sänder en frame och läser tillbaka en dominant ACK-lucka.
\begin{enumerate}[a)]
  \item Vad vet den säkert?
  \item Vad vet den inte?
  \item Noden är ensam på bussen och sänder. Vad läser den i ACK-luckan, och vad ska en
    standardenlig kontroller göra då?
\end{enumerate}

\exercise{DLC}{Resonemang}
\label{ex:2:4}
DLC är fyra bitar men bara värdena 0--8 är meningsfulla.
\begin{enumerate}[a)]
  \item Varför gjordes fältet fyra bitar brett?
  \item En drivrutin skriver DLC 9 till sin kontroller. Vems ansvar är det att det inte händer,
    och varför just den sidan?
\end{enumerate}
